% !TEX spellcheck = en_US



% ~~~~~~~~~~~~~~~~~~~~~~~~~~~~~~~~~~~~~~~ V
% (NC) 2026 Haluk Bingol
% github.com/halukbingol/LaTeX-Templates
%
% Licensees may copy, distribute, display, and perform the work and make derivative works 
% and remixes based on it only for non-commercial purposes. 
% ~~~~~~~~~~~~~~~~~~~~~~~~~~~~~~~~~~~~~~~ A




% =====================================================================
%  Generic Stack and Queue -- Freshman Level
%  ---------------------------------------------------------------------
%  Built on hblib/Template/hblibTemplate.tex with the libraries in ../../hblib/.
%  Programs and their real outputs are in codeoutput/:
%     codeoutput/<Name>.java   a program (classes without main() are only compiled)
%     codeoutput/<Name>.txt    what it printed
% =====================================================================

\documentclass[aspectratio=169,10pt,compress]{beamer}

% ~~~~~~~~~~~~~~~~~~~~~~~~~~~~~~~~~~~~~~~~~~~~~~~~~~~~~~~~~~~~~~~~~ V hblib
	\usepackage{../../hblib/hblibPresentation} % lib presentation: theme, i/N, \hSection, algorithm2e, ...
	\usepackage{../../hblib/hblibCodeoutput} % lib code + output: \lstcode, \lstterminal, ...
	\usepackage{../../hblib/hblibDatastructures} % lib data structures: \hString, \hStack, \hQueue
% ~~~~~~~~~~~~~~~~~~~~~~~~~~~~~~~~~~~~~~~~~~~~~~~~~~~~~~~~~~~~~~~~~ A hblib


% xcolor, array (and the L{width} column), listings, tikz and algorithm2e come from the hblib libraries
\usepackage[T1]{fontenc}
\usepackage{lmodern}
\usepackage{booktabs}

\newcommand{\win}{\textbf{Windows}}
\newcommand{\mac}{\textbf{macOS}}

% listings: fixed character cells, so the columns of a program's output stay aligned
\lstset{columns=fixed}

% algorithm2e set-up (the package is loaded by hblibPresentation)
\SetAlCapFnt{\small}
\SetAlCapNameFnt{\small}
\SetKwProg{Fn}{function}{}{end}
\SetKw{KwError}{error}

% ---------------------------------------------------------------------
\title{Set, Stack and Queue}
\subtitle{Freshman Level: Your First Generic Containers in Java}
\author[Alpaydin]{Ruhan Alpaydin}

\institute[FCIS]{
    Faculty of Computer and Information Sciences\\
    Yeditepe University
}
\date{
	{\tiny \hbTimeStamp}
}

\begin{document}




% ~~~~~~~~~~~~~~~~~~~~~~~~~~~~~~~~~~~~~~ V frm
\begin{frame}[t,fragile]
  \titlepage
\end{frame}
% ~~~~~~~~~~~~~~~~~~~~~~~~~~~~~~~~~~~~~~ A frm




% ~~~~~~~~~~~~~~~~~~~~~~~~~~~~~~~~~~~~~~ V frm
\begin{frame}[t,fragile] \frametitle{Outline}
  \begin{columns}[T]
    \begin{column}{0.48\textwidth}
      \tableofcontents[sections={1-3}, hideallsubsections]
    \end{column}
    \begin{column}{0.48\textwidth}
      \tableofcontents[sections={4-5}, hideallsubsections]
    \end{column}
  \end{columns}
\end{frame}
% ~~~~~~~~~~~~~~~~~~~~~~~~~~~~~~~~~~~~~~ A frm

% ~~~~~~~~~~~~~~~~~~~~~~~~~~~~~~~~~~~~~~ V
\hSection[Ready]{Getting Ready}
% ~~~~~~~~~~~~~~~~~~~~~~~~~~~~~~~~~~~~~~ A




% ~~~~~~~~~~~~~~~~~~~~~~~~~~~~~~~~~~~~~~ V frm
\begin{frame}[t,fragile] \frametitle{What you need}
  \begin{itemize}
    \item A \textbf{JDK} 17 or newer: it has \hCode{javac} (compiler) and \hCode{java} (runner).
    \item A terminal and a text editor (or an IDE such as Eclipse, IntelliJ IDEA).
    \item One folder for this week, e.g.\ \hCode{generic}.
  \end{itemize}
  \begin{columns}[T]
    \begin{column}{0.48\textwidth}
      \begin{block}{Windows (Command Prompt)}
\begin{lstlisting}[style=codesnippet, language={}, basicstyle=\ttfamily\scriptsize]
winget install EclipseAdoptium.Temurin.21.JDK
rem open a NEW Command Prompt, then:
mkdir generic
cd generic
\end{lstlisting}
      \end{block}
    \end{column}
    \begin{column}{0.48\textwidth}
      \begin{block}{macOS (Terminal)}
\begin{lstlisting}[style=codesnippet, language=bash, basicstyle=\ttfamily\scriptsize]
brew install --cask temurin@21
# or the .pkg from adoptium.net
mkdir generic
cd generic
\end{lstlisting}
      \end{block}
    \end{column}
  \end{columns}
  \small Check with:
  \lstterminal[basicstyle=\ttfamily\scriptsize]{codeoutput/Check.txt}
\end{frame}
% ~~~~~~~~~~~~~~~~~~~~~~~~~~~~~~~~~~~~~~ A frm

% ~~~~~~~~~~~~~~~~~~~~~~~~~~~~~~~~~~~~~~ V
\hSection[Set]{The Set}
% ~~~~~~~~~~~~~~~~~~~~~~~~~~~~~~~~~~~~~~ A




% ~~~~~~~~~~~~~~~~~~~~~~~~~~~~~~~~~~~~~~ V frm
\begin{frame}[t,fragile] \frametitle{Set, Stack, Queue: one picture}
  \begin{columns}[T]
    \begin{column}{0.32\textwidth}
      \begin{block}{Set}
        \centering\small
        \hCode{\{apple, banana, cherry\}}\\[0.3em]
        \emph{no duplicates}
      \end{block}
      \small \hRemark{Is this element in the set?}
    \end{column}
    \begin{column}{0.32\textwidth}
      \begin{block}{Stack (LIFO)}
        \begin{center}
        \begin{tikzpicture}\hStack{0}{5}{3}{A,B,C,D}[0.5][0.4]\end{tikzpicture}
        \end{center}
        \small \hCode{pop} $\to$ \textbf{D}
      \end{block}
      \small \hRemark{Undo, brackets, backtracking.}
    \end{column}
    \begin{column}{0.32\textwidth}
      \begin{block}{Queue (FIFO)}
        \begin{center}
        \begin{tikzpicture}\hQueueArray{6}{0}{4}{4}{A,B,C,D}[0.5][0.4]\end{tikzpicture}
        \end{center}
        \small \hCode{dequeue} $\to$ \textbf{A}
      \end{block}
      \small \hRemark{Printers, buffers, BFS.}
    \end{column}
  \end{columns}
  \vspace{0.5em}
  \small All three are \hRemark{generic}: you choose the element type in angle brackets.
\end{frame}
% ~~~~~~~~~~~~~~~~~~~~~~~~~~~~~~~~~~~~~~ A frm




% ~~~~~~~~~~~~~~~~~~~~~~~~~~~~~~~~~~~~~~ V frm
\begin{frame}[t,fragile] \frametitle{What ``generic'' means}
  \begin{columns}[T]
    \begin{column}{0.5\textwidth}
      \begin{itemize}
        \item A \hRemark{generic type} is parameterized by the type of its elements.
        \item \hCode{Set<String>}, \hCode{Stack<Integer>}, \hCode{Queue<Person>}.
        \item The \emph{same class} works for any reference type.
        \item The compiler inserts the casts: \hCode{names.get(0)} is already a \hCode{String}.
      \end{itemize}
    \end{column}
    \begin{column}{0.46\textwidth}
      \begin{block}{Why generics}
        \small
        \begin{itemize}
          \item Type-safe: a wrong type is caught at \textbf{compile time}.
          \item No \hCode{ClassCastException} from forgotten casts.
        \end{itemize}
      \end{block}
    \end{column}
  \end{columns}
  \vspace{0.4em}
  \begin{center}\footnotesize
  \begin{tabular}{@{}lll@{}}
    \toprule
    \textbf{ADT}   & \textbf{Rule}                 & \textbf{Java type} \\
    \midrule
    \textbf{Set}   & no duplicates                 & \hCode{java.util.Set} \\
    \textbf{Stack} & LIFO                          & \hCode{Stack}, \hCode{Deque} \\
    \textbf{Queue} & FIFO                          & \hCode{java.util.Queue} \\
    \bottomrule
  \end{tabular}
  \end{center}
\end{frame}
% ~~~~~~~~~~~~~~~~~~~~~~~~~~~~~~~~~~~~~~ A frm




% ~~~~~~~~~~~~~~~~~~~~~~~~~~~~~~~~~~~~~~ V frm
\begin{frame}[t,fragile] \frametitle{What is a Set?}
  \begin{itemize}
    \item A \hRemark{Set} is a collection of \textbf{unique} elements:
          add the same value twice, and the second add is ignored.
    \item Order is \emph{not} part of the contract, unless the implementation says so.
    \item Uses: keywords, visited cells, tags, membership tests, deduplication.
  \end{itemize}
  \begin{block}{Two properties}
    \begin{itemize}
      \item Does not allow duplicate elements.
      \item Does not maintain insertion order (for \hCode{HashSet}).
    \end{itemize}
  \end{block}
\end{frame}
% ~~~~~~~~~~~~~~~~~~~~~~~~~~~~~~~~~~~~~~ A frm




% ~~~~~~~~~~~~~~~~~~~~~~~~~~~~~~~~~~~~~~ V frm
\begin{frame}[t,fragile] \frametitle{Three Set implementations}
  \begin{center}\small
  \begin{tabular}{L{0.20\textwidth}L{0.26\textwidth}L{0.14\textwidth}L{0.28\textwidth}}
    \toprule
    \textbf{Class} & \textbf{Ordering}             & \textbf{Time}   & \textbf{When to use} \\
    \midrule
    \hCode{HashSet}       & none                          & $O(1)$ avg  & default; fast lookup \\
    \hCode{TreeSet}       & sorted (natural/comparator)   & $O(\log n)$ & sorted iteration \\
    \hCode{LinkedHashSet} & insertion order               & $O(1)$ avg  & predictable order \\
    \bottomrule
  \end{tabular}
  \end{center}
  \vspace{0.3em}
  \begin{itemize}\small
    \item All three implement \hCode{java.util.Set}.
    \item \hCode{HashSet} uses \hCode{hashCode()} and \hCode{equals()}.
    \item \hCode{TreeSet} uses \hCode{compareTo()} (or a \hCode{Comparator}).
  \end{itemize}
\end{frame}
% ~~~~~~~~~~~~~~~~~~~~~~~~~~~~~~~~~~~~~~ A frm




% ~~~~~~~~~~~~~~~~~~~~~~~~~~~~~~~~~~~~~~ V frm
\begin{frame}[t,fragile] \frametitle{Key methods of \texttt{Set}}
  \begin{center}\small
  \begin{tabular}{lL{0.58\textwidth}}
    \toprule
    \textbf{Method}              & \textbf{Meaning} \\
    \midrule
    \hCode{add(E item)}         & add; \hCode{false} if it was already present \\
    \hCode{remove(Object o)}    & remove; \hCode{true} if found \\
    \hCode{contains(Object o)}  & membership test \\
    \hCode{isEmpty()}           & is there no element? \\
    \hCode{size()}              & number of elements \\
    \bottomrule
  \end{tabular}
  \end{center}
  \vspace{0.3em}
  \begin{itemize}\small
    \item \hCode{add} returns \hCode{boolean}: a free ``was it new?'' test.
    \item \hCode{contains} is the workhorse: sets exist for this call.
  \end{itemize}
\end{frame}
% ~~~~~~~~~~~~~~~~~~~~~~~~~~~~~~~~~~~~~~ A frm




% ~~~~~~~~~~~~~~~~~~~~~~~~~~~~~~~~~~~~~~ V frm
\begin{frame}[t,fragile] \frametitle{Set: first example}
  \begin{columns}[T]
    \begin{column}{0.5\textwidth}
      \lstcode[basicstyle=\ttfamily\tiny]{codeoutput/SetDemo.java}{Java}
    \end{column}
    \begin{column}{0.46\textwidth}
      \lstterminal[basicstyle=\ttfamily\scriptsize]{codeoutput/SetDemo.txt}
      \small
      \begin{itemize}
        \item The second \hCode{add("Apple")} returns \hCode{false}.
        \item \hCode{LinkedHashSet} keeps \hCode{[Apple, Banana]}.
        \item \hCode{HashSet} does \emph{not} promise any order.
      \end{itemize}
    \end{column}
  \end{columns}
\end{frame}
% ~~~~~~~~~~~~~~~~~~~~~~~~~~~~~~~~~~~~~~ A frm




% ~~~~~~~~~~~~~~~~~~~~~~~~~~~~~~~~~~~~~~ V frm
\begin{frame}[t,fragile] \frametitle{Set: tracing the calls}
  \begin{center}\small
  \begin{tabular}{lll}
    \toprule
    \textbf{Call}                  & \textbf{Returns} & \textbf{Set contents} \\
    \midrule
    \hCode{add("Apple")}          & \hCode{true}    & \{Apple\} \\
    \hCode{add("Banana")}         & \hCode{true}    & \{Apple, Banana\} \\
    \hCode{add("Apple")}          & \hCode{false}   & \{Apple, Banana\} \\
    \hCode{contains("Banana")}    & \hCode{true}    & \{Apple, Banana\} \\
    \hCode{size()}                & \hCode{2}       & \{Apple, Banana\} \\
    \hCode{remove("Apple")}       & \hCode{true}    & \{Banana\} \\
    \hCode{isEmpty()}             & \hCode{false}   & \{Banana\} \\
    \bottomrule
  \end{tabular}
  \end{center}
  \vspace{0.3em}
  \small Nothing in this table ever has two copies of the same value:
  that is the whole contract.
\end{frame}
% ~~~~~~~~~~~~~~~~~~~~~~~~~~~~~~~~~~~~~~ A frm




% ~~~~~~~~~~~~~~~~~~~~~~~~~~~~~~~~~~~~~~ V frm
\begin{frame}[t,fragile] \frametitle{\texttt{HashSet} vs \texttt{LinkedHashSet}}
  \begin{columns}[T]
    \begin{column}{0.56\textwidth}
      \lstcode[basicstyle=\ttfamily\tiny]{codeoutput/SetOrder.java}{Java}
    \end{column}
    \begin{column}{0.4\textwidth}
      \lstterminal[basicstyle=\ttfamily\scriptsize]{codeoutput/SetOrder.txt}
      \small
      Same input, \hRemark{different iteration order}.
      Pick \hCode{LinkedHashSet} when you need the output to be predictable.
    \end{column}
  \end{columns}
\end{frame}
% ~~~~~~~~~~~~~~~~~~~~~~~~~~~~~~~~~~~~~~ A frm




% ~~~~~~~~~~~~~~~~~~~~~~~~~~~~~~~~~~~~~~ V frm
\begin{frame}[t,fragile] \frametitle{Exercise 1: remove duplicates using a Set}
  \begin{block}{Problem}
    Write a Java program to remove duplicates from a list using a \hCode{Set}.
  \end{block}
  \begin{columns}[T]
    \begin{column}{0.54\textwidth}
      \lstcode[basicstyle=\ttfamily\tiny]{codeoutput/RemoveDup.java}{Java}
    \end{column}
    \begin{column}{0.42\textwidth}
      \lstterminal[basicstyle=\ttfamily\scriptsize]{codeoutput/RemoveDup.txt}
      \small
      One line does it all:
      \hCode{new LinkedHashSet<>(words)}.
      \hCode{HashSet} is faster but loses the order.
    \end{column}
  \end{columns}
\end{frame}
% ~~~~~~~~~~~~~~~~~~~~~~~~~~~~~~~~~~~~~~ A frm




% ~~~~~~~~~~~~~~~~~~~~~~~~~~~~~~~~~~~~~~ V frm
\begin{frame}[t,fragile] \frametitle{Exercise 2: common elements in two lists}
  \begin{block}{Problem}
    Write a Java program to find common elements between two lists using a \hCode{Set}.
  \end{block}
  \begin{columns}[T]
    \begin{column}{0.54\textwidth}
      \lstcode[basicstyle=\ttfamily\tiny]{codeoutput/CommonElements.java}{Java}
    \end{column}
    \begin{column}{0.42\textwidth}
      \lstterminal[basicstyle=\ttfamily\scriptsize]{codeoutput/CommonElements.txt}
      \small
      \hCode{retainAll} keeps only the elements also present in the
      argument: the \textbf{intersection} of the two sets.
    \end{column}
  \end{columns}
\end{frame}
% ~~~~~~~~~~~~~~~~~~~~~~~~~~~~~~~~~~~~~~ A frm

% ~~~~~~~~~~~~~~~~~~~~~~~~~~~~~~~~~~~~~~ V
\hSection[Stack]{The Stack}
% ~~~~~~~~~~~~~~~~~~~~~~~~~~~~~~~~~~~~~~ A




% ~~~~~~~~~~~~~~~~~~~~~~~~~~~~~~~~~~~~~~ V frm
\begin{frame}[t,fragile] \frametitle{What is a Stack?}
  \begin{columns}[T]
    \begin{column}{0.48\textwidth}
      \begin{itemize}
        \item Insert and remove at the \textbf{top}.
        \item The most recent element leaves first: \hRemark{LIFO}.
        \item Uses: undo (Ctrl+Z), brackets, expression evaluation,
              backtracking, DFS.
      \end{itemize}
    \end{column}
    \begin{column}{0.48\textwidth}
      \begin{center}
      \begin{tikzpicture}\hStack{0}{6}{3}{A,B,C,D}[0.55][0.45]\end{tikzpicture}
      \end{center}
      \small \hCode{pop} returns \textbf{D}.
    \end{column}
  \end{columns}
\end{frame}
% ~~~~~~~~~~~~~~~~~~~~~~~~~~~~~~~~~~~~~~ A frm




% ~~~~~~~~~~~~~~~~~~~~~~~~~~~~~~~~~~~~~~ V frm
\begin{frame}[t,fragile] \frametitle{LIFO: last in, first out}
  \begin{columns}[T]
    \begin{column}{0.5\textwidth}
      \begin{center}\footnotesize
      \begin{tabular}{lll}
        \toprule
        \textbf{Call}             & \textbf{Returns} & \textbf{Stack (top $\to$ bottom)} \\
        \midrule
        \hCode{push("A")}        &                  & A \\
        \hCode{push("B")}        &                  & B A \\
        \hCode{push("C")}        &                  & C B A \\
        \hCode{peek()}           & C                & C B A \\
        \hCode{pop()}            & C                & B A \\
        \hCode{push("D")}        &                  & D B A \\
        \hCode{pop()}            & D                & B A \\
        \bottomrule
      \end{tabular}
      \end{center}
    \end{column}
    \begin{column}{0.46\textwidth}
      \begin{block}{LIFO}
        The element added last is the first to come out:\\
        \hRemark{Last In, First Out}.
      \end{block}
      \small Compare with a queue: first in, first out. Same elements,
      same arrival order, \hRemark{opposite order of leaving}.
    \end{column}
  \end{columns}
\end{frame}
% ~~~~~~~~~~~~~~~~~~~~~~~~~~~~~~~~~~~~~~ A frm




% ~~~~~~~~~~~~~~~~~~~~~~~~~~~~~~~~~~~~~~ V frm
\begin{frame}[t,fragile] \frametitle{Class: \texttt{java.util.Stack}}
  \begin{center}\small
  \begin{tabular}{lL{0.55\textwidth}}
    \toprule
    \textbf{Method} & \textbf{Meaning} \\
    \midrule
    \hCode{push(E item)}   & add on top \\
    \hCode{pop()}          & remove and return the top; error if empty \\
    \hCode{peek()}         & return the top without removing it \\
    \hCode{isEmpty()}      & is there no element? \\
    \bottomrule
  \end{tabular}
  \end{center}
  \vspace{0.3em}
  \begin{itemize}\small
    \item \hCode{Stack<E>} is a legacy class (Java 1.0) that extends
          \hCode{Vector}: every call is \hCode{synchronized}.
    \item Modern Java prefers \hCode{Deque<E>} (\hCode{ArrayDeque})
          with the same \hCode{push}/\hCode{pop}/\hCode{peek} names.
  \end{itemize}
\end{frame}
% ~~~~~~~~~~~~~~~~~~~~~~~~~~~~~~~~~~~~~~ A frm




% ~~~~~~~~~~~~~~~~~~~~~~~~~~~~~~~~~~~~~~ V frm
\begin{frame}[t,fragile] \frametitle{Stack: example}
  \begin{columns}[T]
    \begin{column}{0.52\textwidth}
      \lstcode[basicstyle=\ttfamily\tiny]{codeoutput/StackDemo.java}{Java}
    \end{column}
    \begin{column}{0.44\textwidth}
      \lstterminal[basicstyle=\ttfamily\scriptsize]{codeoutput/StackDemo.txt}
      \small After the two \hCode{pop}s the stack is empty: a further
      \hCode{pop} would throw \hCode{EmptyStackException}.
    \end{column}
  \end{columns}
\end{frame}
% ~~~~~~~~~~~~~~~~~~~~~~~~~~~~~~~~~~~~~~ A frm




% ~~~~~~~~~~~~~~~~~~~~~~~~~~~~~~~~~~~~~~ V frm
\begin{frame}[t,fragile] \frametitle{Stack: tracing the calls}
  \begin{center}\small
  \begin{tabular}{lll}
    \toprule
    \textbf{Call}                & \textbf{Returns} & \textbf{Stack (top $\to$ bottom)} \\
    \midrule
    \hCode{push("Java")}        &                  & Java \\
    \hCode{push("Python")}      &                  & Python, Java \\
    \hCode{peek()}              & Python           & Python, Java \\
    \hCode{pop()}               & Python           & Java \\
    \hCode{pop()}               & Java             & (empty) \\
    \hCode{isEmpty()}           & \hCode{true}     & (empty) \\
    \bottomrule
  \end{tabular}
  \end{center}
  \vspace{0.3em}
  \small \hCode{peek} does not change the stack; \hCode{pop} removes the
  top and returns it.
\end{frame}
% ~~~~~~~~~~~~~~~~~~~~~~~~~~~~~~~~~~~~~~ A frm




% ~~~~~~~~~~~~~~~~~~~~~~~~~~~~~~~~~~~~~~ V frm
\begin{frame}[t,fragile] \frametitle{\texttt{pop} on an empty stack}
  \begin{columns}[T]
    \begin{column}{0.5\textwidth}
      \lstcode[basicstyle=\ttfamily\tiny]{codeoutput/StackBug.java}{Java}
    \end{column}
    \begin{column}{0.46\textwidth}
      \lstterminal[basicstyle=\ttfamily\scriptsize]{codeoutput/StackBug.txt}
      \small
      \begin{itemize}
        \item \hCode{ArrayDeque.pop} throws \hCode{NoSuchElementException} on empty.
        \item Fix: guard with \hCode{!isEmpty()}, or use \hCode{poll()}
              (returns \hCode{null}).
      \end{itemize}
    \end{column}
  \end{columns}
\end{frame}
% ~~~~~~~~~~~~~~~~~~~~~~~~~~~~~~~~~~~~~~ A frm




% ~~~~~~~~~~~~~~~~~~~~~~~~~~~~~~~~~~~~~~ V frm
\begin{frame}[t,fragile] \frametitle{Modern alternative: \texttt{Deque}}
  \begin{columns}[T]
    \begin{column}{0.52\textwidth}
      \lstcode[basicstyle=\ttfamily\tiny]{codeoutput/DequeStack.java}{Java}
    \end{column}
    \begin{column}{0.44\textwidth}
      \lstterminal[basicstyle=\ttfamily\scriptsize]{codeoutput/DequeStack.txt}
      \small
      \begin{itemize}
        \item Same \hCode{push}/\hCode{pop}/\hCode{peek} names as \hCode{Stack}.
        \item \hCode{ArrayDeque} is faster and forbids \hCode{null}.
        \item \hRemark{Rule of thumb}: use \hCode{ArrayDeque} for both
              stack and queue.
      \end{itemize}
    \end{column}
  \end{columns}
\end{frame}
% ~~~~~~~~~~~~~~~~~~~~~~~~~~~~~~~~~~~~~~ A frm




% ~~~~~~~~~~~~~~~~~~~~~~~~~~~~~~~~~~~~~~ V frm
\begin{frame}[t,fragile] \frametitle{Exercise 3: reverse a string using a Stack}
  \begin{block}{Problem}
    Write a Java program to reverse a string using a \hCode{Stack}.
  \end{block}
  \begin{columns}[T]
    \begin{column}{0.54\textwidth}
      \lstcode[basicstyle=\ttfamily\tiny]{codeoutput/ReverseString.java}{Java}
    \end{column}
    \begin{column}{0.42\textwidth}
      \lstterminal[basicstyle=\ttfamily\scriptsize]{codeoutput/ReverseString.txt}
      \small
      Push every character, then pop them: the last pushed is the
      first to come out, so the string is reversed.
    \end{column}
  \end{columns}
\end{frame}
% ~~~~~~~~~~~~~~~~~~~~~~~~~~~~~~~~~~~~~~ A frm

% ~~~~~~~~~~~~~~~~~~~~~~~~~~~~~~~~~~~~~~ V
\hSection[Queue]{The Queue}
% ~~~~~~~~~~~~~~~~~~~~~~~~~~~~~~~~~~~~~~ A




% ~~~~~~~~~~~~~~~~~~~~~~~~~~~~~~~~~~~~~~ V frm
\begin{frame}[t,fragile] \frametitle{What is a Queue?}
  \begin{columns}[T]
    \begin{column}{0.48\textwidth}
      \begin{itemize}
        \item Insert at the \textbf{rear}, remove at the \textbf{front}.
        \item The oldest element leaves first: \hRemark{FIFO}.
        \item Uses: printers, keyboard buffers, scheduling, BFS,
              message queues.
      \end{itemize}
    \end{column}
    \begin{column}{0.48\textwidth}
      \begin{center}
      \begin{tikzpicture}\hQueueArray{6}{0}{4}{4}{A,B,C,D}[0.55][0.45]\end{tikzpicture}
      \end{center}
      \small \hCode{dequeue} returns \textbf{A}.
    \end{column}
  \end{columns}
\end{frame}
% ~~~~~~~~~~~~~~~~~~~~~~~~~~~~~~~~~~~~~~ A frm




% ~~~~~~~~~~~~~~~~~~~~~~~~~~~~~~~~~~~~~~ V frm
\begin{frame}[t,fragile] \frametitle{FIFO: first in, first out}
  \begin{columns}[T]
    \begin{column}{0.5\textwidth}
      \begin{center}\footnotesize
      \begin{tabular}{lll}
        \toprule
        \textbf{Call}           & \textbf{Returns} & \textbf{Queue (front $\to$ rear)} \\
        \midrule
        \hCode{offer(A)}       &                  & A \\
        \hCode{offer(B)}       &                  & A B \\
        \hCode{offer(C)}       &                  & A B C \\
        \hCode{poll()}         & A                & B C \\
        \hCode{offer(D)}       &                  & B C D \\
        \hCode{peek()}         & B                & B C D \\
        \hCode{poll()}         & B                & C D \\
        \bottomrule
      \end{tabular}
      \end{center}
    \end{column}
    \begin{column}{0.46\textwidth}
      \begin{block}{FIFO}
        The element that has waited the longest is served first:\\
        \hRemark{First In, First Out}.
      \end{block}
      \small Same four elements, same arrival order, \hRemark{opposite
      order of leaving} from a stack.
    \end{column}
  \end{columns}
\end{frame}
% ~~~~~~~~~~~~~~~~~~~~~~~~~~~~~~~~~~~~~~ A frm




% ~~~~~~~~~~~~~~~~~~~~~~~~~~~~~~~~~~~~~~ V frm
\begin{frame}[t,fragile] \frametitle{\texttt{java.util.Queue}: two forms}
  \begin{itemize}
    \item \hCode{Queue} is an \textbf{interface}, not a class.
    \item Main implementations: \hCode{ArrayDeque}, \hCode{LinkedList}.
    \item Every operation has two forms: one \textbf{throws},
          one returns a \textbf{special value}.
  \end{itemize}
  \begin{center}\small
  \begin{tabular}{lll}
    \toprule
    \textbf{Operation} & \textbf{throws exception} & \textbf{returns special value} \\
    \midrule
    insert    & \hCode{add(e)}     & \hCode{offer(e)} \\
    remove    & \hCode{remove()}   & \hCode{poll()} \\
    examine   & \hCode{element()}  & \hCode{peek()} \\
    \bottomrule
  \end{tabular}
  \end{center}
  \small \hCode{isEmpty()} and \hCode{size()} complete the interface.
\end{frame}
% ~~~~~~~~~~~~~~~~~~~~~~~~~~~~~~~~~~~~~~ A frm




% ~~~~~~~~~~~~~~~~~~~~~~~~~~~~~~~~~~~~~~ V frm
\begin{frame}[t,fragile] \frametitle{Queue: example}
  \begin{columns}[T]
    \begin{column}{0.52\textwidth}
      \lstcode[basicstyle=\ttfamily\tiny]{codeoutput/QueueDemo.java}{Java}
    \end{column}
    \begin{column}{0.44\textwidth}
      \lstterminal[basicstyle=\ttfamily\scriptsize]{codeoutput/QueueDemo.txt}
      \small \hCode{ArrayDeque} prints from the \textbf{front};
      \hCode{poll()} on empty returns \hCode{null}.
    \end{column}
  \end{columns}
\end{frame}
% ~~~~~~~~~~~~~~~~~~~~~~~~~~~~~~~~~~~~~~ A frm




% ~~~~~~~~~~~~~~~~~~~~~~~~~~~~~~~~~~~~~~ V frm
\begin{frame}[t,fragile] \frametitle{\texttt{PriorityQueue} is not FIFO}
  \begin{columns}[T]
    \begin{column}{0.52\textwidth}
      \lstcode[basicstyle=\ttfamily\tiny]{codeoutput/PriorityDemo.java}{Java}
    \end{column}
    \begin{column}{0.44\textwidth}
      \lstterminal[basicstyle=\ttfamily\scriptsize]{codeoutput/PriorityDemo.txt}
      \small
      It implements \hCode{Queue}, but the element that leaves is the
      \textbf{smallest} (or per a \hCode{Comparator}), not the oldest.
    \end{column}
  \end{columns}
\end{frame}
% ~~~~~~~~~~~~~~~~~~~~~~~~~~~~~~~~~~~~~~ A frm




% ~~~~~~~~~~~~~~~~~~~~~~~~~~~~~~~~~~~~~~ V frm
\begin{frame}[t,fragile] \frametitle{Exercise 4: queue from two stacks}
  \begin{block}{Problem}
    Write a Java program to implement a queue using two stacks.
  \end{block}
  \begin{columns}[T]
    \begin{column}{0.54\textwidth}
      \lstcode[basicstyle=\ttfamily\tiny]{codeoutput/QueueTwoStacks.java}{Java}
    \end{column}
    \begin{column}{0.42\textwidth}
      \lstterminal[basicstyle=\ttfamily\scriptsize]{codeoutput/QueueTwoStacks.txt}
      \small
      Two reversals compose to the original order:
      \hCode{in} receives, \hCode{out} serves; when \hCode{out} is
      empty, transfer everything.
    \end{column}
  \end{columns}
\end{frame}
% ~~~~~~~~~~~~~~~~~~~~~~~~~~~~~~~~~~~~~~ A frm

% ~~~~~~~~~~~~~~~~~~~~~~~~~~~~~~~~~~~~~~ V
\hSection[Wrap]{Wrap-up}
% ~~~~~~~~~~~~~~~~~~~~~~~~~~~~~~~~~~~~~~ A




% ~~~~~~~~~~~~~~~~~~~~~~~~~~~~~~~~~~~~~~ V frm
\begin{frame}[t,fragile] \frametitle{The errors you will see first}
  \lstterminal[basicstyle=\ttfamily\scriptsize]{codeoutput/Typos.txt}
  \vspace{0.2em}
  \small
  Fix the \emph{first} error first. Java is \textbf{case sensitive}:
  \hCode{SetDemo} is not \hCode{setdemo}, and \hCode{java trace}
  is not \hCode{java Trace}.
\end{frame}
% ~~~~~~~~~~~~~~~~~~~~~~~~~~~~~~~~~~~~~~ A frm




% ~~~~~~~~~~~~~~~~~~~~~~~~~~~~~~~~~~~~~~ V frm
\begin{frame}[t,fragile] \frametitle{Cheat sheet}
  \begin{columns}[T]
    \begin{column}{0.54\textwidth}
      \begin{center}\scriptsize
      \begin{tabular}{@{}ll@{}}
        \toprule
        \textbf{Task}             & \textbf{Code} \\
        \midrule
        unique elements           & \hCode{new HashSet<>()} \\
        sorted set                & \hCode{new TreeSet<>()} \\
        insertion-ordered set     & \hCode{new LinkedHashSet<>()} \\
        \midrule
        LIFO container            & \hCode{Deque<E> s = new ArrayDeque<>()} \\
        stack push                & \hCode{s.push(x)} \\
        stack pop                 & \hCode{s.pop()} \\
        stack peek                & \hCode{s.peek()} \\
        \midrule
        FIFO container            & \hCode{Queue<E> q = new ArrayDeque<>()} \\
        queue enqueue             & \hCode{q.offer(x)} \\
        queue dequeue             & \hCode{q.poll()} \\
        queue peek                & \hCode{q.peek()} \\
        \bottomrule
      \end{tabular}
      \end{center}
    \end{column}
    \begin{column}{0.42\textwidth}
      \begin{block}{Do}
        \scriptsize
        use \hCode{ArrayDeque} for both stack and queue;
        prefer \hCode{offer}/\hCode{poll}/\hCode{peek} when the
        container may be empty; guard \hCode{pop} with \hCode{isEmpty()}.
      \end{block}
      \begin{block}{Don't}
        \scriptsize
        rely on \hCode{HashSet} order; mix \hCode{addLast} with
        \hCode{removeLast} and call it a queue; store \hCode{null}
        in an \hCode{ArrayDeque}; treat \hCode{PriorityQueue} as FIFO.
      \end{block}
    \end{column}
  \end{columns}
\end{frame}
% ~~~~~~~~~~~~~~~~~~~~~~~~~~~~~~~~~~~~~~ A frm




% ~~~~~~~~~~~~~~~~~~~~~~~~~~~~~~~~~~~~~~ V frm
\begin{frame}[t,fragile] \frametitle{Practice, and questions}
  \begin{enumerate}\small
    \item Read a sentence and print how many \textbf{distinct} words it has.
    \item Read a word, print its reverse using a \hCode{Deque} as a stack.
    \item Read numbers until \texttt{0}, print them in the same order with a \hCode{Queue}.
    \item Use a \hCode{Set<Integer>} to print every distinct value of an \hCode{int[]}.
    \item Build a small \emph{undo} with a stack: type \texttt{undo} to remove the last word.
    \item Fix \hCode{StackBug}: guard the second \hCode{pop()} with \hCode{!isEmpty()}.
  \end{enumerate}
  \vspace{0.5em}
  \begin{center}
    \hRemark{\large Check yourself: does \hCode{add} tell you it was new?
    Does \hCode{poll} throw on empty? Is \hCode{PriorityQueue} FIFO?}
  \end{center}
  \vspace{0.6em}
  \begin{center}
    \hRemark{\Huge Questions?}
  \end{center}
\end{frame}
% ~~~~~~~~~~~~~~~~~~~~~~~~~~~~~~~~~~~~~~ A frm




\end{document}